% =====================================================================
%  Bien dich: pdfLaTeX (mac dinh cua Overleaf) -- khong can chinh gi.
%  Neu muon dung font he thong thi doi compiler sang XeLaTeX va thay
%  ba dong inputenc/fontenc/lmodern ben duoi bang:
%      \usepackage{fontspec}
%      \setmainfont{TeX Gyre Termes}
% =====================================================================
\documentclass[11pt,a4paper]{article}

\usepackage[utf8]{inputenc}
\usepackage[T5]{fontenc}
\usepackage{lmodern}

\usepackage[margin=2.6cm]{geometry}
\usepackage{amsmath}
\usepackage{amssymb}
\usepackage{booktabs}
\usepackage{array}
\usepackage{caption}
\usepackage{microtype}
\usepackage{enumitem}
\usepackage{tikz}
\usetikzlibrary{arrows.meta}
\usepackage[hidelinks]{hyperref}

\renewcommand{\tablename}{Bảng}
\renewcommand{\figurename}{Hình}
\renewcommand{\abstractname}{Tóm tắt}
\renewcommand{\refname}{Tài liệu tham khảo}
\renewcommand{\contentsname}{Nội dung}
\renewcommand{\appendixname}{Phụ lục}

\captionsetup[table]{skip=6pt}
\captionsetup[figure]{skip=8pt}
\newcolumntype{R}{>{\raggedleft\arraybackslash}p{2.7cm}}

\title{\bfseries Hợp nhất hai tầng cho HRM-PET:\\
định tuyến đa nguồn và phân lớp bổ sung với cơ chế cổng theo biên}
\author{Nguyễn Thiện Trường}
\date{Tháng 8, 2026}

\begin{document}
\maketitle

\begin{abstract}
Chúng tôi thay bộ định tuyến đơn nguồn của HRM-PET bằng một phép hợp nhất
(fusion) hai nguồn tín hiệu, sau đó sử dụng chính nguồn thứ hai làm bộ phân lớp
bổ sung với trọng số được điều khiển bởi một cơ chế cổng (gating) theo từng mẫu.
Nguồn thứ hai là một đầu phân lớp dựa trên phép chiếu ngẫu nhiên bậc hai theo
kiểu RanPAC, xây dựng trực tiếp trên đặc trưng LoRA mà HRM-PET đã huấn luyện.
Nhờ vậy phương pháp không cần huấn luyện thêm tham số nào và vẫn tuân thủ giao
thức không lưu mẫu (exemplar-free). Trên ba bộ dữ liệu học liên tục và bốn seed,
phương pháp cải thiện 17 trong 18 phép so sánh: độ chính xác định tuyến tăng từ
$+1.10$ đến $+2.26$ điểm, độ chính xác phân lớp tăng từ $+0.69$ đến $+1.50$
điểm. Chi phí tính toán tăng thêm đúng bằng một lượt truyền xuôi (forward pass)
trên mỗi mẫu.
\end{abstract}

\tableofcontents
\bigskip

\section{Giới thiệu}

\subsection{Quy trình của HRM-PET}

HRM-PET xử lý mỗi mẫu đầu vào qua bốn bước:

\begin{enumerate}[leftmargin=*, itemsep=2pt]
\item \textbf{Suy luận danh tính tác vụ} (task-identity inference, TII): ước
      lượng mẫu thuộc về tác vụ nào, dựa trên năng lượng prompt.
\item \textbf{Định tuyến}: chuyển mẫu tới bộ tham số LoRA đã đóng băng của tác
      vụ được chọn ở bước 1.
\item \textbf{Tinh chỉnh}: hai cơ chế DRM và CRM điều chỉnh lại kết quả định
      tuyến.
\item \textbf{Phân lớp}: một đầu phân lớp dùng chung cho mọi tác vụ đưa ra dự
      đoán cuối cùng.
\end{enumerate}

Toàn bộ chất lượng của chuỗi xử lý này phụ thuộc vào một tín hiệu định tuyến duy
nhất, sinh ra ở bước 1.

\subsection{Hai hạn chế}

Báo cáo chỉ ra hai hạn chế của thiết kế trên. Thứ nhất, tín hiệu TII không phải
tín hiệu định tuyến tốt nhất có thể khai thác từ chính mô hình đã huấn luyện.
Thứ hai, ngay cả khi định tuyến được cải thiện, phần lớn mức cải thiện đó không
chuyển hóa thành độ chính xác phân lớp, vì lý do nằm ở kiến trúc chứ không ở
chất lượng định tuyến. Mục \ref{sec:analysis} trình bày số liệu cho cả hai.

\subsection{Đóng góp}

\begin{enumerate}[leftmargin=*, itemsep=3pt]
\item \textbf{Hợp nhất ở cấp định tuyến.} Kết hợp TII với một đầu phân lớp dựa
      trên phép chiếu ngẫu nhiên, sau khi chuẩn hóa hai nguồn về cùng thang đo.
      Hai nguồn cho kết quả sai trên những tập mẫu khác nhau, nên việc kết hợp
      đem lại lợi ích thực (Mục \ref{sec:stage1}).
\item \textbf{Hợp nhất ở cấp phân lớp với cơ chế cổng theo biên.} Sử dụng đầu
      chiếu ngẫu nhiên như một bộ phân lớp thứ hai, với trọng số tỉ lệ nghịch
      với mức độ chắc chắn của đầu phân lớp chính. Cơ chế cổng loại bỏ hoàn toàn
      suy giảm hàm mất mát mà phiên bản dùng trọng số cố định gây ra
      (Mục \ref{sec:stage2}).
\item \textbf{Một trục kiểm chứng chính xác.} Tham số $w$ của tầng thứ nhất nội
      suy liên tục về HRM-PET gốc; tại $w = 1.0$, phương pháp tái lập baseline
      đến từng chữ số (Mục \ref{sec:sanity}).
\end{enumerate}

Cả hai thành phần đều đặt ở giai đoạn suy luận nên không yêu cầu huấn luyện lại
mô hình. Hình \ref{fig:pipeline} minh họa vị trí của chúng trong quy trình.

\begin{figure}[htbp]
\centering
\resizebox{\textwidth}{!}{%
\begin{tikzpicture}[
  font=\footnotesize,
  >={Stealth[length=2mm]},
  box/.style={draw, minimum height=8mm, minimum width=1.6cm, inner xsep=4pt, align=center},
  hl/.style={box, very thick, fill=black!6}
]
\node[box] (in) at (0.0, 0.0)    {Ảnh vào};
\node[box] (tii) at (2.7, 0.85)  {TII};
\node[box] (rp) at (2.7, -0.85)  {Đầu RP};
\node[hl]  (f1) at (5.9, 0.0)    {Tầng 1\\Hợp nhất định tuyến};
\node[box] (lo) at (8.9, 0.0)    {Chọn LoRA};
\node[box] (dr) at (11.3, 0.0)   {DRM / CRM};
\node[box] (cl) at (13.9, 0.0)   {Đầu phân lớp};
\node[hl]  (f2) at (13.9, -2.1)  {Tầng 2\\Hợp nhất phân lớp};
\node[box] (out) at (13.9, -3.7) {Dự đoán};

\draw[->] (in) -- (tii);
\draw[->] (in) -- (rp);
\draw[->] (tii) -- (f1);
\draw[->] (rp) -- (f1);
\draw[->] (f1) -- (lo);
\draw[->] (lo) -- (dr);
\draw[->] (dr) -- (cl);
\draw[->] (cl) -- (f2);
\draw[->] (f2) -- (out);
\draw[->, dashed] (rp.south) |- node[pos=0.75, above, font=\scriptsize]
      {điểm số theo lớp của đầu RP} (f2.west);
\end{tikzpicture}}
\caption{Quy trình suy luận. Các khối viền đậm là hai thành phần được bổ sung.
Đầu RP được tính một lần duy nhất và cung cấp tín hiệu cho cả hai tầng: điểm số
theo tác vụ cho tầng thứ nhất, điểm số theo lớp cho tầng thứ hai (đường nét
đứt).}
\label{fig:pipeline}
\end{figure}

\section{Phân tích hai hạn chế}
\label{sec:analysis}

\subsection{Hạn chế thứ nhất: tín hiệu định tuyến chưa tối ưu}

Chúng tôi xây dựng một đầu phân lớp dựa trên phép chiếu ngẫu nhiên (random
projection, viết tắt là đầu RP) trên chính đặc trưng LoRA của HRM-PET. Đầu RP
cho độ chính xác định tuyến cao hơn TII ở cả ba bộ dữ liệu. Quan trọng hơn,
\emph{hợp} của hai bộ định tuyến vượt xa kết quả mà HRM-PET đạt được sau khi đã
tinh chỉnh bằng DRM và CRM.

\begin{table}[htbp]
\centering
\caption{Độ chính xác định tuyến (\%), seed 42. Hàng cuối biểu thị tỉ lệ mẫu mà
ít nhất một trong hai bộ định tuyến cho kết quả đúng, tức là giới hạn trên có
thể đạt được khi kết hợp hai nguồn.}
\label{tab:routing}
\begin{tabular}{lrrr}
\toprule
Nguồn định tuyến & ImageNet-R & CIFAR-100 & CUB-200 \\
\midrule
TII                        & 63.80 & 81.21 & 92.96 \\
HRM-PET sau DRM và CRM     & 77.79 & 89.82 & 93.21 \\
Đầu RP                     & 77.12 & 89.51 & 94.02 \\
\midrule
Hợp của TII và đầu RP      & \textbf{82.99} & \textbf{92.86} & \textbf{96.26} \\
\bottomrule
\end{tabular}
\end{table}

Hai bộ định tuyến cho kết quả sai trên những tập mẫu khác nhau. Trên ImageNet-R,
đầu RP định tuyến đúng $19.2\,\%$ số mẫu mà TII định tuyến sai. Tính bổ trợ này
là điều kiện cần để việc hợp nhất có ý nghĩa: nếu hai nguồn cùng sai trên một
tập mẫu, hợp của chúng sẽ không cao hơn nguồn mạnh hơn.

\subsection{Hạn chế thứ hai: định tuyến tốt hơn không đồng nghĩa phân lớp tốt hơn}

Sau khi cải thiện tầng định tuyến, một hiện tượng khác xuất hiện: mức tăng của
độ chính xác định tuyến chỉ chuyển hóa thành độ chính xác phân lớp với tỉ lệ
$14$--$21\,\%$, và trên CUB-200 thì hoàn toàn không chuyển hóa.

Nguyên nhân nằm ở kiến trúc. Đầu phân lớp dùng chung trải trên \emph{toàn bộ}
các lớp đã quan sát, nên một mẫu bị định tuyến sai vẫn thường được phân lớp
đúng. Trong trường hợp đó, việc sửa lại đường định tuyến không đem lại lợi ích
nào; ngược lại, định tuyến lại một mẫu vốn đã được phân lớp đúng có thể làm kết
quả trở nên sai.

Tuy nhiên, đầu RP không chỉ là một bộ định tuyến. Bản thân nó là một bộ phân lớp
hoàn chỉnh mà tầng định tuyến đã bỏ qua. Bảng \ref{tab:cls} đo mức bổ trợ giữa
hai bộ phân lớp.

\begin{table}[htbp]
\centering
\caption{Tính bổ trợ ở cấp phân lớp (\%), seed 42. Cột cuối biểu thị tỉ lệ mẫu
mà đầu phân lớp sau định tuyến cho kết quả sai còn đầu RP cho kết quả đúng.}
\label{tab:cls}
\begin{tabular}{lrrrr}
\toprule
Bộ dữ liệu & Đầu sau định tuyến & Đầu RP & Hợp & Phần đầu RP bổ trợ \\
\midrule
ImageNet-R & 74.31 & 70.08 & 78.89 & $+4.58$ \\
CIFAR-100  & 89.97 & 88.76 & 92.40 & $+2.43$ \\
CUB-200    & 86.48 & 87.50 & 90.35 & $+3.87$ \\
\bottomrule
\end{tabular}
\end{table}

Ngay trên ImageNet-R, nơi đầu RP có độ chính xác tổng thể thấp hơn $4$ điểm, nó
vẫn phân lớp đúng $4.58\,\%$ số mẫu mà đầu phân lớp chính không xử lý được.

\section{Phương pháp đề xuất}

\subsection{Nguồn tín hiệu thứ hai}

Nguồn tín hiệu thứ hai gồm một phép chiếu ngẫu nhiên đóng băng, sinh lại được từ
một seed cố định, theo sau là hàm phi tuyến ReLU và một lời giải hồi quy ridge
trên ma trận Gram tích lũy. Không có tham số nào được huấn luyện. Đặc trưng $f$
được trích từ \emph{một} bộ tham số LoRA cố định của HRM-PET, đóng vai trò thích
nghi phiên đầu (first-session adaptation):
\begin{align}
h &= \mathrm{ReLU}(f W), \qquad W \text{ đóng băng, sinh lại từ seed},\\
G &= \sum_n h_n h_n^{\top}, \qquad
C = \sum_n h_n\,\mathrm{onehot}(y_n)^{\top},\\
s &= h^{\top} (G + \lambda I)^{-1} C .
\end{align}

Hệ thống chỉ lưu hai đại lượng tổng hợp $G$ và $C$; ma trận chiếu $W$ được sinh
lại từ seed. Không có đặc trưng của từng mẫu hay ảnh của các tác vụ trước được
lưu trữ, nên giao thức không lưu mẫu vẫn được bảo toàn.

\subsection{Tầng thứ nhất: hợp nhất ở cấp định tuyến}
\label{sec:stage1}

Ký hiệu $z(\cdot)$ là phép chuẩn hóa về trung bình $0$ và phương sai $1$, tính
riêng cho từng mẫu trên tập các tác vụ đã quan sát. Điểm số theo tác vụ của hai
nguồn được kết hợp tuyến tính; tác vụ có điểm cao nhất được chọn, và kết quả
định tuyến này được đưa vào các cơ chế DRM, CRM cùng đầu phân lớp sẵn có của
HRM-PET:
\begin{equation}
\mathrm{fused} = w\, z(\ell^{\mathrm{TII}}) + (1-w)\, z(s^{\mathrm{RP}}),
\qquad w = 0.7 .
\label{eq:route}
\end{equation}

Phép chuẩn hóa ở đây là bắt buộc. Hàm $\log\mathrm{softmax}$ bảo toàn độ phân
tán của đầu vào, trong khi logit của TII có độ phân tán lớn hơn nhiều so với
điểm số hồi quy ridge. Nếu kết hợp trực tiếp, TII sẽ chi phối kết quả ở mọi giá
trị trọng số, và độ chính xác định tuyến giảm xuống \emph{thấp hơn} cả mức của
riêng đầu RP: $68.24$ so với $77.12$.

\subsection{Tầng thứ hai: hợp nhất ở cấp phân lớp với cơ chế cổng}
\label{sec:stage2}

Phiên bản đầu tiên kết hợp hai bộ phân lớp bằng một trọng số $\beta$ cố định cho
mọi mẫu. Cách làm này gây ra một suy giảm có tính hệ thống: giá trị hàm mất mát
entropy chéo (cross-entropy) trên CIFAR-100 tăng $+0.019 \pm 0.003$, dù độ chính
xác vẫn được cải thiện.

Nguyên nhân là trọng số cố định xử lý như nhau hai loại mẫu có bản chất khác
hẳn: mẫu mà đầu phân lớp chính dự đoán với xác suất $p = 0.97$, và mẫu mà nó
chưa phân định rõ giữa các lớp. Trên CIFAR-100, loại thứ nhất chiếm đa số. Việc
dành $30\,\%$ trọng số quyết định cho nguồn thứ hai trên chính những mẫu đó chỉ
làm phẳng phân bố xác suất tại một dự đoán vốn đã đúng: cross-entropy tăng ngay,
trong khi kết quả $\arg\max$ không thay đổi.

Giải pháp là cho $\beta$ thay đổi theo từng mẫu, tỉ lệ nghịch với mức độ chắc
chắn của đầu phân lớp chính. Đại lượng dùng để đo mức độ chắc chắn là biên
(margin) giữa hai lớp có điểm cao nhất của chính nó. Cơ sở của lựa chọn này:
việc hợp nhất chỉ có thể thay đổi dự đoán bằng cách thay thế lớp xếp thứ nhất,
mà ứng viên thay thế khả dĩ nhất là lớp xếp thứ hai. Gọi $p_{(1)}$ và $p_{(2)}$
là hai xác suất lớn nhất của đầu phân lớp sau định tuyến:
\begin{equation}
\beta_i = \beta \left( 1 - \frac{p_{(1)} - p_{(2)}}{p_{(1)} + p_{(2)}} \right),
\qquad \beta = 0.5 .
\label{eq:gate}
\end{equation}
Điểm số cuối cùng được kết hợp rồi đưa trở lại thang đo của logit ban đầu:
\begin{align}
m_i &= (1 - \beta_i)\, z(\ell_i) + \beta_i\, z(s_i),\\
\hat{\ell}_i &= m_i \cdot \mathrm{std}(\ell_i) + \mathrm{mean}(\ell_i).
\label{eq:affine}
\end{align}

Hai ví dụ minh họa tác dụng của cơ chế cổng. Với mẫu có $p_{(1)} = 0.90$ và
$p_{(2)} = 0.02$, trọng số hiệu dụng $\beta_i \approx 0.007$, tức là dự đoán gần
như được giữ nguyên. Với mẫu có $p_{(1)} = 0.35$ và $p_{(2)} = 0.30$, trọng số
$\beta_i \approx 0.28$ và nguồn thứ hai đóng góp gần như tối đa. Cơ chế cổng
không đưa thêm siêu tham số nào và áp dụng một quy tắc duy nhất cho cả ba bộ dữ
liệu.

Phép biến đổi \eqref{eq:affine} là biến đổi affine theo từng mẫu, do đó đơn
điệu: thứ tự xếp hạng giữa các lớp được bảo toàn và mọi chỉ số độ chính xác
không thay đổi. Vai trò duy nhất của nó là đưa thang đo của logit về mức so sánh
được với baseline khi tính cross-entropy.

\section{Thiết lập thí nghiệm}

\subsection{Cấu hình}

Mạng xương sống (backbone) là ViT-B/16 tiền huấn luyện trên ImageNet-21K, thích
nghi bằng LoRA với hạng $8$. Mỗi bộ dữ liệu được chia thành chuỗi $10$ tác vụ và
thí nghiệm được lặp trên bốn seed $42$, $43$, $44$, $45$. Toàn bộ checkpoint của
LoRA và TII được huấn luyện lại từ đầu cho từng seed. Khâu đánh giá sử dụng cùng
một bộ checkpoint cho cả baseline lẫn phương pháp đề xuất, nên chênh lệch quan
sát được phản ánh đúng phần thay đổi ở giai đoạn suy luận, không lẫn nhiễu do
huấn luyện.

Siêu tham số của đầu RP: số chiều chiếu $10^4$, hàm phi tuyến ReLU, hệ số điều
chuẩn $\lambda = 10^4$. Không áp dụng hiệu chỉnh nhiệt độ (temperature scaling).

\subsection{Các chỉ số đánh giá}

Ký hiệu $a_{i,t}$ là độ chính xác trên tác vụ $i$ đo tại giai đoạn $t$, và $T$
là tổng số tác vụ.

\begin{itemize}[leftmargin=*, itemsep=2pt]
\item \textbf{Acc@task}: độ chính xác định tuyến, tức tỉ lệ mẫu được gán đúng
      tác vụ.
\item \textbf{Acc@1}, \textbf{Acc@5}: độ chính xác phân lớp ở nhãn xếp thứ nhất
      và trong năm nhãn dẫn đầu.
\item \textbf{Loss}: giá trị hàm mất mát cross-entropy.
\item \textbf{Forgetting} và \textbf{Backward}: hai chỉ số đo mức độ mô hình giữ
      lại kiến thức cũ, tính trên toàn bộ ma trận kết quả.
\end{itemize}

\begin{equation}
\mathrm{Forgetting} = \frac{1}{T}\sum_{i=1}^{T}
\Big( \max_{t} a_{i,t} - a_{i,T} \Big),
\qquad
\mathrm{Backward} = \frac{1}{T}\sum_{i=1}^{T} \big( a_{i,T} - a_{i,i} \big).
\end{equation}

Với Loss và Forgetting, giá trị càng thấp càng tốt; với bốn chỉ số còn lại, giá
trị càng cao càng tốt. Điểm cần lưu ý: cả Forgetting và Backward đều so sánh một
tác vụ với chính nó ở thời điểm trước đó, chứ không so sánh giữa các phương
pháp. Đặc điểm này sẽ được dùng lại ở Mục \ref{sec:ramp}.

\section{Kết quả}

Bảng \ref{tab:main} tổng hợp kết quả trên cả ba bộ dữ liệu. Giá trị là trung
bình $\pm$ độ lệch chuẩn trên bốn seed. Cột \emph{Thay đổi} được tính theo cặp:
chênh lệch giữa phương pháp đề xuất và baseline được tính riêng cho từng seed,
sau đó mới lấy trung bình. Nhờ vậy, độ lệch chuẩn của cột này phản ánh độ ổn
định của chính mức cải thiện, không bị ảnh hưởng bởi độ khó khác nhau giữa các
seed.

\begin{table}[htbp]
\centering
\caption{Kết quả trên ba bộ dữ liệu, trung bình $\pm$ độ lệch chuẩn trên bốn
seed 42, 43, 44, 45.}
\label{tab:main}
\begin{tabular}{@{}p{0.4cm}lRRR@{}}
\toprule
\multicolumn{2}{@{}l}{Bộ dữ liệu và chỉ số} & HRM-PET & Đề xuất & Thay đổi \\
\midrule
\multicolumn{5}{l}{\emph{Split-ImageNet-R}}\\
& Acc@task   & $77.58 \pm 0.33$  & $79.85 \pm 0.18$  & $+2.264 \pm 0.158$ \\
& Acc@1      & $73.94 \pm 0.48$  & $75.18 \pm 0.24$  & $+1.236 \pm 0.416$ \\
& Acc@5      & $86.71 \pm 0.21$  & $88.02 \pm 0.17$  & $+1.312 \pm 0.148$ \\
& Loss       & $1.23 \pm 0.01$   & $1.19 \pm 0.01$   & $-0.040 \pm 0.002$ \\
& Forgetting & $3.36 \pm 0.56$   & $3.30 \pm 0.35$   & $-0.056 \pm 0.334$ \\
& Backward   & $-3.16 \pm 0.61$  & $-3.20 \pm 0.42$  & $-0.047 \pm 0.282$ \\
\addlinespace
\multicolumn{5}{l}{\emph{Split-CIFAR100}}\\
& Acc@task   & $89.77 \pm 0.11$  & $90.87 \pm 0.09$  & $+1.103 \pm 0.066$ \\
& Acc@1      & $89.78 \pm 0.03$  & $90.47 \pm 0.05$  & $+0.690 \pm 0.063$ \\
& Acc@5      & $98.65 \pm 0.03$  & $98.86 \pm 0.07$  & $+0.217 \pm 0.057$ \\
& Loss       & $0.39 \pm 0.00$   & $0.39 \pm 0.00$   & $-0.000 \pm 0.001$ \\
& Forgetting & $3.89 \pm 0.08$   & $3.74 \pm 0.12$   & $-0.156 \pm 0.098$ \\
& Backward   & $-3.87 \pm 0.07$  & $-3.73 \pm 0.11$  & $+0.142 \pm 0.097$ \\
\addlinespace
\multicolumn{5}{l}{\emph{Split-CUB200}}\\
& Acc@task   & $93.12 \pm 0.14$  & $94.34 \pm 0.10$  & $+1.226 \pm 0.149$ \\
& Acc@1      & $86.38 \pm 0.25$  & $87.88 \pm 0.27$  & $+1.495 \pm 0.353$ \\
& Acc@5      & $97.04 \pm 0.10$  & $97.44 \pm 0.13$  & $+0.400 \pm 0.103$ \\
& Loss       & $0.57 \pm 0.00$   & $0.53 \pm 0.01$   & $-0.036 \pm 0.004$ \\
& Forgetting & $2.54 \pm 0.35$   & $2.18 \pm 0.12$   & $-0.359 \pm 0.325$ \\
& Backward   & $-2.42 \pm 0.39$  & $-2.14 \pm 0.13$  & $+0.279 \pm 0.391$ \\
\bottomrule
\end{tabular}
\end{table}

Ba nhận xét về bảng trên.

\textbf{Mức độ bao phủ.} Trong tổng số $18$ phép so sánh ($6$ chỉ số trên $3$ bộ
dữ liệu), $17$ phép cho kết quả cải thiện. Trường hợp còn lại là Backward trên
ImageNet-R với giá trị $-0.047 \pm 0.282$: độ lệch chuẩn lớn gấp sáu lần trị
tuyệt đối của trung bình, nên đại lượng này không phân biệt được với $0$. Đây là
dao động giữa các seed, không phải suy giảm có tính hệ thống.

\textbf{Mức độ tin cậy.} Hai chỉ số đáng tin cậy nhất là Acc@task và Acc@5, với
độ lệch chuẩn của mức thay đổi nhỏ hơn chính mức thay đổi khoảng mười đến mười
lăm lần. Nói cách khác, mức cải thiện ở hai chỉ số này lặp lại ổn định trên mọi
seed.

\textbf{Trường hợp CUB-200.} Bộ dữ liệu này vốn là điểm yếu của phiên bản chỉ có
tầng định tuyến: mức tăng độ chính xác định tuyến không chuyển hóa thành Acc@1
mà dừng ở $-0.03$. Sau khi bổ sung tầng phân lớp, đây trở thành bộ dữ liệu có
Acc@1 tăng mạnh nhất, đúng như phân tích ở Mục \ref{sec:analysis}.

\section{Kiểm chứng tính đúng đắn của cài đặt}
\label{sec:sanity}

Đặt $w = 1.0$ trong \eqref{eq:route} sẽ dồn toàn bộ trọng số về TII, tương đương
với việc vô hiệu hóa nguồn tín hiệu thứ hai. Ở cấu hình này, phương pháp tái lập
chính xác baseline HRM-PET đến từng chữ số trên cả ba bộ dữ liệu: $77.7914$,
$74.0191$, $1.2305$ và các giá trị còn lại, không sai khác tại chữ số thập phân
thứ tư.

Kết quả này có hai hệ quả. Thứ nhất, mọi chênh lệch báo cáo ở Bảng \ref{tab:main}
đến từ đúng các thành phần được bổ sung, không phải từ khác biệt ngẫu nhiên
trong quy trình đánh giá. Thứ hai, $w$ là một tham số liên tục nội suy giữa
phương pháp đề xuất và baseline, do đó cung cấp sẵn một trục khảo sát loại bỏ
thành phần (ablation) liên tục thay vì chỉ một so sánh nhị phân.

\section{Khảo sát loại bỏ thành phần}

\subsection{Trọng số hợp nhất ở cấp phân lớp}

\begin{table}[htbp]
\centering
\caption{Ảnh hưởng của $\beta$ và của cơ chế cổng. Hai hàng đầu đo trên ba seed,
hàng cuối đo trên seed 42.}
\label{tab:beta}
\begin{tabular}{lrrl}
\toprule
Cấu hình & Loss trên CIFAR-100 & Backward trên ImageNet-R & Nhận xét \\
\midrule
$\beta = 0.3$, không có cổng  & $+0.019 \pm 0.003$ & $-0.298 \pm 0.391$ & suy giảm có hệ thống \\
$\beta = 0.5$, cổng theo biên & $-0.000 \pm 0.001$ & $-0.047 \pm 0.282$ & cấu hình được chọn \\
$\beta = 0.8$, cổng theo biên & $+0.029$           & $-0.462$           & hợp nhất quá mức \\
\bottomrule
\end{tabular}
\end{table}

Ở cấu hình $\beta = 0.3$, độ lệch chuẩn của Loss chỉ bằng một phần sáu giá trị
trung bình, cho thấy đây là suy giảm thực và lặp lại trên mọi seed. Đặc điểm này
khác hẳn chỉ số Backward, nơi độ lệch chuẩn lớn hơn trung bình nhiều lần và do
đó không cho phép kết luận về dấu.

Do cơ chế cổng làm giảm đáng kể trọng số hiệu dụng, giá trị $\beta$ danh nghĩa
cần tăng từ $0.3$ lên $0.5$. Khi tăng tiếp lên $0.8$, mức hợp nhất trở nên quá
lớn và ImageNet-R suy giảm đồng thời ở Loss, Forgetting và Backward. Giá trị tối
ưu do đó nằm trong khoảng giữa chứ không ở biên của dải khảo sát.

\subsection{Trọng số tăng dần theo giai đoạn: đã cài đặt, đã đo và đã loại bỏ}
\label{sec:ramp}

Như đã lưu ý ở phần định nghĩa chỉ số, cả Forgetting và Backward đều so sánh một
tác vụ với đỉnh hiệu năng của chính nó. Hệ quả là một hiệu chỉnh áp dụng đồng
đều ở mọi giai đoạn sẽ triệt tiêu khỏi cả hai chỉ số. Điều này giải thích tại
sao hai chỉ số đó gần như không thay đổi trong Bảng \ref{tab:main}, đồng thời
gợi ý cho trọng số hợp nhất tăng dần theo số tác vụ đã quan sát, theo dạng
$(\text{số tác vụ đã quan sát}/T)^{\gamma}$.

Cơ chế này hoạt động đúng như thiết kế. Với $\gamma = 2$, Forgetting cải thiện
lên $-0.415$, $-0.433$, $-0.462$ và Backward lên $+0.357$, $+0.411$, $+0.462$
trên ba bộ dữ liệu, trong khi dòng kết quả tại giai đoạn cuối không thay đổi một
chữ số nào, do hệ số nhân đạt giá trị $1.0$ tại giai đoạn đó.

Chính đặc điểm sau cùng lại là lý do để loại bỏ cơ chế này. Nếu kết quả tại giai
đoạn cuối không đổi mà hai chỉ số so với đỉnh vẫn cải thiện, thì mức cải thiện
đó chỉ có thể đến từ việc hiệu năng của mô hình giảm ở các giai đoạn trung gian.
Độ chính xác trung bình tính trên cả mười giai đoạn xác nhận điều này: trên
CUB-200, giá trị giảm từ $88.35$ xuống $88.16$ rồi $87.91$. Cải thiện chỉ số ghi
nhớ bằng cách hạ thấp độ chính xác ở các giai đoạn trung gian không phải là một
cải tiến thực chất.

Hai kết quả có giá trị vẫn thu được từ thí nghiệm này:

\begin{enumerate}[leftmargin=*, itemsep=3pt]
\item Cơ chế hợp nhất ở cấp phân lớp \emph{thực sự} gây hại tại giai đoạn đầu.
      Trên ImageNet-R, độ chính xác ở tác vụ thứ nhất là $86.0$ so với $87.5$
      của baseline, phù hợp với lập luận rằng ước lượng ma trận Gram cần đủ số
      lớp mới đáng tin cậy.
\item Cũng chính đầu RP đó lại là bộ định tuyến mạnh ngay từ giai đoạn đầu. Do
      vậy, việc giảm trọng số định tuyến ở các giai đoạn sớm chỉ làm giảm độ
      chính xác định tuyến trung bình trên cả ba bộ dữ liệu mà không đem lại lợi
      ích bù lại.
\end{enumerate}

\section{Chi phí tính toán}

\begin{table}[htbp]
\centering
\caption{Số lượt gọi LoRA trung bình trên mỗi mẫu.}
\label{tab:cost}
\begin{tabular}{lrrr}
\toprule
Bộ dữ liệu & HRM-PET & Đề xuất & Thay đổi \\
\midrule
ImageNet-R & 1.83 & 2.85 & $+1.02$ \\
CIFAR-100  & 1.40 & 2.41 & $+1.01$ \\
CUB-200    & 1.60 & 2.60 & $+1.00$ \\
\bottomrule
\end{tabular}
\end{table}

Chi phí tăng thêm đúng bằng một lượt truyền xuôi cho đầu RP. Tầng hợp nhất ở cấp
phân lớp không phát sinh chi phí bổ sung, vì nó sử dụng lại chính điểm số đã
được tính cho tầng định tuyến, như minh họa bằng đường nét đứt trong Hình
\ref{fig:pipeline}.

Chi phí bộ nhớ tăng thêm gồm một ma trận Gram, một ma trận tương quan lớp và một
seed để sinh lại ma trận chiếu.

\section{Hạn chế}

\textbf{Chỉ số Backward trên ImageNet-R.} Giá trị $-0.047 \pm 0.282$ không phân
biệt được với $0$. Cần thêm số lượng seed để kết luận về dấu, tuy nhiên có thể
khẳng định đây không phải suy giảm có tính hệ thống.

\textbf{Giới hạn cải thiện của Loss.} Đây là hệ quả số học chứ không phải thiếu
sót của phương pháp. Với cross-entropy quanh mức $1.23$, mức tăng $1.24$ điểm
của Acc@1 chỉ tương ứng khoảng $0.01$--$0.04$ nats, đúng bằng mức đã đạt được.
Muốn cải thiện thêm thì phải áp dụng hiệu chỉnh nhiệt độ, nhưng hiệu chỉnh cho
phương pháp đề xuất trong khi giữ nguyên baseline là một so sánh không công
bằng; nếu hiệu chỉnh cho cả hai thì chênh lệch trở lại như cũ.

\textbf{Cách chọn siêu tham số.} Giá trị $\beta$ được chọn trên seed 42 và dùng
chung cho cả ba bộ dữ liệu, sau đó được kiểm chứng lại trên ba seed còn lại.
Chúng tôi không đặt siêu tham số riêng cho từng bộ dữ liệu.

\textbf{Phạm vi đánh giá.} Thí nghiệm được thực hiện trên ba bộ dữ liệu. Mã
nguồn đã có sẵn bộ nạp dữ liệu cho ImageNet-A, nhưng dữ liệu cần được tải thủ
công nên bộ này chưa được đưa vào đánh giá.

\section{Kết luận}

Hai tầng hợp nhất, cùng đặt ở giai đoạn suy luận và không yêu cầu huấn luyện
lại, cải thiện $17$ trong $18$ phép so sánh trên ba bộ dữ liệu và bốn seed. Tầng
định tuyến khai thác tính bổ trợ giữa hai nguồn tín hiệu vốn cho kết quả sai
trên những tập mẫu khác nhau. Tầng phân lớp khai thác thực tế rằng nguồn thứ hai
bản thân nó đã là một bộ phân lớp hoàn chỉnh nhưng bị kiến trúc gốc bỏ qua.

Cơ chế cổng theo biên là thành phần giúp tầng thứ hai không phải trả giá trên
những mẫu mà mô hình gốc vốn đã dự đoán chắc chắn. Cơ chế này loại bỏ suy giảm
có tính hệ thống duy nhất mà chúng tôi quan sát được, và không đưa thêm siêu
tham số nào vào phương pháp.

\appendix

\section{Đối chiếu thuật ngữ}

\begin{table}[htbp]
\centering
\caption{Các thuật ngữ tiếng Anh và cách dùng tương ứng trong báo cáo.}
\label{tab:glossary}
\begin{tabular}{ll}
\toprule
Tiếng Anh & Dùng trong báo cáo \\
\midrule
continual learning              & học liên tục \\
exemplar-free                   & không lưu mẫu \\
task-identity inference (TII)   & suy luận danh tính tác vụ \\
routing / router                & định tuyến / bộ định tuyến \\
fusion                          & hợp nhất \\
gating                          & cơ chế cổng \\
margin                          & biên \\
random projection               & phép chiếu ngẫu nhiên \\
frozen                          & đóng băng \\
first-session adaptation        & thích nghi phiên đầu \\
backbone                        & mạng xương sống \\
forward pass                    & lượt truyền xuôi \\
cross-entropy                   & entropy chéo \\
temperature scaling             & hiệu chỉnh nhiệt độ \\
ablation study                  & khảo sát loại bỏ thành phần \\
complementarity                 & tính bổ trợ \\
baseline                        & baseline \\
seed                            & seed \\
checkpoint                      & checkpoint \\
\bottomrule
\end{tabular}
\end{table}

% Truoc khi nop: bo sung trich dan cho HRM-PET va RanPAC.

\end{document}
